% ========================================================
% BÀI TẬP TOÁN 5 - KẾT NỐI TRI THỨC VỚI CUỘC SỐNG
% Bài 67: Luyện tập chung
% Định dạng: BTTeX Part III - Trắc nghiệm điền số ngắn
% ========================================================

% Câu 1
\begin{ex}
  Thống kê điểm số của 8 bạn: 7, 8, 9, 8, 10, 8, 9, 5. Số học sinh đạt điểm 8 là bao nhiêu bạn?
  \shortans{3}
  \loigiai{
    Có 3 bạn điểm 8.
  }
\end{ex}

% Câu 2
\begin{ex}
  Ở bảng trên, hiệu giữa điểm số lớn nhất và nhỏ nhất là bao nhiêu?
  \shortans{5}
  \loigiai{
    $10 - 5 = 5$.
  }
\end{ex}

% Câu 3
\begin{ex}
  Biểu đồ quạt tròn biểu diễn 3 môn thể thao: Bóng đá $50\%$, Cầu lông $30\%$, Bơi lội $20\%$. Có 120 học sinh tham gia. Số bạn thích Bơi lội là bao nhiêu bạn?
  \shortans{24}
  \loigiai{
    $120 \times 20\% = 24$ bạn.
  }
\end{ex}

% Câu 4
\begin{ex}
  Ở trường trên, số bạn thích Bóng đá là bao nhiêu bạn?
  \shortans{60}
  \loigiai{
    $120 \times 50\% = 60$ bạn.
  }
\end{ex}

% Câu 5
\begin{ex}
  Ở trường trên, số bạn thích Cầu lông là bao nhiêu bạn?
  \shortans{36}
  \loigiai{
    $120 \times 30\% = 36$ bạn.
  }
\end{ex}

% Câu 6
\begin{ex}
  Tung đồng xu 50 lần, có 28 lần mặt sấp. Số lần mặt ngửa xuất hiện là bao nhiêu lần?
  \shortans{22}
  \loigiai{
    $50 - 28 = 22$ lần.
  }
\end{ex}

% Câu 7
\begin{ex}
  Ở thực nghiệm trên, tỉ số phần trăm số lần xuất hiện mặt ngửa là bao nhiêu phần trăm?
  \shortans{44}
  \loigiai{
    $(22 : 50) \times 100\% = 44\%$.
  }
\end{ex}

% Câu 8
\begin{ex}
  Gieo một con xúc xắc 25 lần, mặt 1 chấm xuất hiện 5 lần. Tỉ số phần trăm xuất hiện mặt 1 chấm là bao nhiêu phần trăm?
  \shortans{20}
  \loigiai{
    $(5 : 25) \times 100\% = 20\%$.
  }
\end{ex}

% Câu 9
\begin{ex}
  Trong một hộp có 10 viên bi đỏ và 15 viên bi xanh. Tỉ số số bi xanh so với tổng số bi khi rút gọn có mẫu số là bao nhiêu?
  \shortans{5}
  \loigiai{
    $\frac{15}{25} = \frac{3}{5}$. Mẫu số là 5.
  }
\end{ex}

% Câu 10
\begin{ex}
  Ở hộp bi trên, tỉ số phần trăm số bi đỏ so với tổng số bi là bao nhiêu phần trăm?
  \shortans{40}
  \loigiai{
    $(10 : 25) \times 100\% = 40\%$.
  }
\end{ex}

% Câu 11
\begin{ex}
  Một bảng số liệu có 5 giá trị: 12, 14, 16, 18, 20. Trung bình cộng của 5 giá trị này là bao nhiêu?
  \shortans{16}
  \loigiai{
    $(12 + 14 + 16 + 18 + 20) : 5 = 80 : 5 = 16$.
  }
\end{ex}

% Câu 12
\begin{ex}
  Bảng phân phối điểm kiểm tra có 2 bạn điểm 10, 5 bạn điểm 9, 3 bạn điểm 8. Tổng số bạn tham gia kiểm tra là bao nhiêu?
  \shortans{10}
  \loigiai{
    $2 + 5 + 3 = 10$ bạn.
  }
\end{ex}

% Câu 13
\begin{ex}
  Biểu đồ quạt tròn chia làm 5 phần bằng nhau. Mỗi phần tương ứng bao nhiêu phần trăm?
  \shortans{20}
  \loigiai{
    $100\% : 5 = 20\%$.
  }
\end{ex}

% Câu 14
\begin{ex}
  Một xạ thủ bắn 20 phát đạn, có 18 phát trúng mục tiêu. Tỉ số phần trăm bắn trúng là bao nhiêu phần trăm?
  \shortans{90}
  \loigiai{
    $(18 : 20) \times 100\% = 90\%$.
  }
\end{ex}

% Câu 15
\begin{ex}
  Ở bài xạ thủ trên, có bao nhiêu phát trượt mục tiêu?
  \shortans{2}
  \loigiai{
    $20 - 18 = 2$ phát.
  }
\end{ex}

% Câu 16
\begin{ex}
  Điều tra mức độ hài lòng của 200 khách hàng: Rất hài lòng $40\%$, Hài lòng $50\%$, Không hài lòng $10\%$. Số người rất hài lòng là bao nhiêu người?
  \shortans{80}
  \loigiai{
    $200 \times 40\% = 80$ người.
  }
\end{ex}

% Câu 17
\begin{ex}
  Ở cuộc khảo sát trên, số người không hài lòng là bao nhiêu người?
  \shortans{20}
  \loigiai{
    $200 \times 10\% = 20$ người.
  }
\end{ex}

% Câu 18
\begin{ex}
  Lấy ngẫu nhiên một quả bóng trong hộp 40 lần, có 16 lần bóng màu vàng. Tỉ số phần trăm bóng màu vàng là bao nhiêu phần trăm?
  \shortans{40}
  \loigiai{
    $(16 : 40) \times 100\% = 40\%$.
  }
\end{ex}

% Câu 19
\begin{ex}
  Chi tiêu ăn uống của một gia đình chiếm $50\%$ tổng thu nhập. Nếu tháng đó dành $10$ triệu đồng cho ăn uống thì tổng thu nhập là bao nhiêu triệu đồng?
  \shortans{20}
  \loigiai{
    $10 : 50\% = 20$ triệu đồng.
  }
\end{ex}

% Câu 20
\begin{ex}
  Gieo xúc xắc 100 lần có 48 lần ra mặt số chẵn. Tỉ số phần trăm ra mặt số lẻ là bao nhiêu phần trăm?
  \shortans{52}
  \loigiai{
    $100 - 48 = 52\%$.
  }
\end{ex}

% Câu 21
\begin{ex}
  Một trường tiểu học có 500 học sinh. Biểu đồ quạt cho biết tỉ lệ học sinh nữ là $52\%$. Số học sinh nữ là bao nhiêu bạn?
  \shortans{260}
  \loigiai{
    $500 \times 52\% = 260$ bạn.
  }
\end{ex}

% Câu 22
\begin{ex}
  Ở trường trên, số học sinh nam là bao nhiêu bạn?
  \shortans{240}
  \loigiai{
    $500 - 260 = 240$ bạn.
  }
\end{ex}

% Câu 23
\begin{ex}
  Số bạn nam ít hơn số bạn nữ bao nhiêu bạn?
  \shortans{20}
  \loigiai{
    $260 - 240 = 20$ bạn.
  }
\end{ex}

% Câu 24
\begin{ex}
  Một tổ công nhân có 4 người làm được số sản phẩm: 25, 30, 35, 30. Trung bình mỗi người làm được bao nhiêu sản phẩm?
  \shortans{30}
  \loigiai{
    $(25 + 30 + 35 + 30) : 4 = 120 : 4 = 30$ sản phẩm.
  }
\end{ex}

% Câu 25
\begin{ex}
  Rút một con bài trong bộ bài tây 50 lần, có 12 lần rút được quân cơ. Tỉ số phần trăm số lần rút được quân cơ là bao nhiêu phần trăm?
  \shortans{24}
  \loigiai{
    $(12 : 50) \times 100\% = 24\%$.
  }
\end{ex}
